\section{Experiment Execution and Resource Infrastructure}
\label{sec:execution}

\begin{designbox}
Route every provider call through one gateway, price it, and charge it against an
atomic budget reservation before the work runs --- so that spend is bounded,
attributable, and impossible for a role to escalate on its own. Long experiments
run as supervised background jobs; scarce hardware is coordinated through an
explicit lease.
\end{designbox}

\subsection{One gateway for every call}

All backend execution flows through a single application-level gateway,
\code{RunExecGateway.execute} (\srcpath{core/run\_gateway.py}), which assigns a
\code{call\_id}, propagates the resume-thread identity, stamps timing, and returns
a normalized \code{RunnerResult}. Because there is exactly one interception point,
tracing, metrics, cost accounting, and policy apply uniformly regardless of role or
provider (Codex, Claude Code, Copilot CLI); every provider interaction carries an
id that ties it to the event tape, which is what makes the per-call audit of
Section~\ref{sec:evidence} possible.

\subsection{Pricing and the usage ledger}

Token usage is priced by a small pricing module (\srcpath{core/pricing.py})
that maps a model to input/output rates (default model \code{gpt-5.5}) and quotes
token usage into dollars, including the conversion of Copilot premium-request
counts into USD. Priced usage is accumulated in a usage ledger
(\srcpath{core/usage.py}) as token counts, premium requests, a nano-AIU credit
total, and cost. The ledger is \emph{presence-aware}: it distinguishes a genuine
zero from a provider that omitted usage (Section~\ref{sec:roles}), so no cost is
ever fabricated from missing data.

\subsection{Budget reservations and spend fences}

Spend control is reservation-based and event-backed. Before a call runs, the
runtime creates a budget reservation; the outcome is recorded as one of
\code{budget.reservation.created}, \code{budget.reservation.denied},
\code{budget.reservation.settled}, or \code{budget.reservation.released}, with
\code{budget.unpriced.blocked} and \code{budget.fence\_breach.blocked} for the
cases where a call cannot be priced or would breach a fence
(\srcpath{core/event\_catalog.py}). Crucially, the provider-side spend fences on a
call (\code{max\_budget\_usd}, \code{max\_ai\_credits}) are injected from the
reservation, not from role configuration --- a role cannot raise its own ceiling.
Table~\ref{tab:budget} lists the principal budget defaults, all verified against
the live configuration.

\begin{table}[h]
\centering
\small
\begin{tabularx}{\linewidth}{@{}p{6.4cm} p{2.2cm} X@{}}
\toprule
\textbf{Control} & \textbf{Default} & \textbf{Effect} \\
\midrule
Per-mission preflight cap (\code{per\_mission\_cap\_usd}) & \$30 & max spend admitted for one mission \\
Daily cap (\code{daily\_cap\_usd})           & \$180 & max spend per local day \\
Global daily cap (\code{global\_daily\_cap\_usd}) & 0 (off) & cross-project daily cap; disabled unless set \\
Autonomous missions per interactive run (\code{max\_missions}) & 6 & bound on one supervisor run (the daemon runs continuously) \\
Planner task-iteration budget (\code{..\_BUDGET\_USD}) & \$30 & per-task iteration budget surfaced to the web layer \\
Planner task-iteration cycles (\code{..\_MAX\_CYCLES}) & 6 & max planner iterations per task \\
Max active daemons (\code{ARGUS\_SKILL\_MAX\_ACTIVE\_DAEMONS}) & 2 & concurrency bound across projects \\
\bottomrule
\end{tabularx}
\caption{Principal budget/quota defaults (\code{life/supervisor/\_config.py},
\code{daemon/config.py}, \code{apps/cli/\_core.py}). All are operator-overridable;
the values shown are the live code defaults.}
\label{tab:budget}
\end{table}

\subsection{Provider quota and fencing}

Beyond dollars, a provider may impose its own quota. A per-day Codex quota state is
tracked in a lock-guarded state file (\srcpath{core/provider\_quota.py}); a call
first acquires a \code{ProviderPermit}, and a denial carries a
\code{stop\_kind} such as \code{budget\_exhausted} so the mission ends in a paused
status rather than an error (Section~\ref{sec:runtime}). Related fencing and a
Copilot-specific guard (\srcpath{core/provider\_fencing.py},
\srcpath{core/copilot\_guard.py}) keep a single provider outage from cascading
into a spin.

\subsection{Background jobs and teams}

Not all work is short. When the Engineer starts a long experiment --- for example a
GPU training run --- it launches a supervised background job that has its own
monitor: the monitor checks health on an interval, early-stops on failure, and
reports a terminal result to the operator inbox. The Engineer reads a registry of
in-flight jobs (\srcpath{engineer/background\_subagents.py}) and classifies each as
self-watched or needs-attention; the cadence-wait mechanism that lets it yield to a
job's own monitor instead of busy-polling is described in
Section~\ref{sec:reliability}. A separate team subsystem
(\srcpath{argus\_skill/team/}, roughly 12 modules) supports parallel subagents with
a Curator maintaining the shared skill pool; it is an optional mode, not part of
the single-mission path.

\subsection{Coordinating scarce hardware: the GPU lease}

On managed GPU hosts a scheduler reclaims cards that sit idle, so an operator runs
a low-duty ``keep-alive'' loader to hold them. Argus ships an explicit
\emph{lease} tool for this (\srcpath{argus\_skill/tools/gpu\_lease.py}, with a
standalone loader \srcpath{argus\_skill/tools/gpu\_load.py}): the recommended path,
\code{gpu\_lease run -\/- <cmd>}, frees the cards, records a lease for the job's
lifetime, runs the command, and on exit re-parks the keep-alive only if no other
lease is active. A \code{--detach} mode lets a supervisor own the lease
independently of the mission, and a watchdog re-parks the cards only once they are
unheld, lease-free, and idle. The keep-alive's argv and cwd are recorded so the
tool restarts the exact program the operator started, and process termination only
ever targets the configured match token.

\begin{callout}[Scope of the GPU lease]
The GPU lease is a real, robust \emph{operator/agent-invoked} capability
(\srcpath{argus\_skill/tools/gpu\_lease.py}, \srcpath{tools/setup.py}), not an
automatic runtime service: nothing in the mission scheduler or engineer loop calls
it implicitly. It is presented here as resource-governance tooling the agent may
choose to use, which is what the live integration supports --- no more.
\end{callout}
